
\documentclass[11pt]{article}
\usepackage[margin=1in]{geometry}
\usepackage{amsmath,amssymb,booktabs,hyperref,graphicx,parskip,enumitem}
\usepackage{xcolor}
\hypersetup{colorlinks=true,linkcolor=blue!50!black,urlcolor=blue!50!black}
\setlist{nosep}

\title{Replication Report: (title not recorded)\\[2pt]\large OSTI 2587225 \textbar{} Coverage 7/10, Agreement 8/10}
\author{Rick Stevens \and Ollie (AI Assistant)\\\small Argonne National Laboratory}
\date{2026-04-27}

\begin{document}\maketitle

\section{Source paper}
\begin{itemize}
\item \textbf{Title:} (title not recorded)
\item \textbf{Authors:} Not recorded
\item \textbf{Year:} n/a
\item \textbf{Domain:} Not recorded
\item \textbf{Identifier:} OSTI 2587225
\end{itemize}


\section{Replication objective}
The central claim of the paper concerns not recorded. Our objective was to reproduce the headline numerical/qualitative results within the constraints of the AI-driven Replication Project (24--72\,h, commodity GPU/CPU compute, no author contact). A replication plan is archived at \texttt{2587225-ScaWL-Scaling-k-WL-Weisfeiler-Lehman-Algorithms-in/replication\_plan.pdf}.

\section{Methodology}
We followed the standard Replication Project workflow: ingest the source PDF, extract method and data pointers, draft a replication plan, attempt the smallest end-to-end pipeline that produces a comparable number, then scale up where compute and data permit. Tools used in this replication are catalogued in the meta-paper's computational tools inventory (see \texttt{COMPUTATIONAL\_TOOLS\_INVENTORY.md}).

This paper went through the Tier-Lift pipeline; supplementary notes at \texttt{.openclaw/workspace/24h-progress/tier-lift-v2/2587225.md}.

\subsection*{Paper-directory README excerpt}
\small
\par\medskip\textbf{ScaWL Replication (OSTI 2587225)}\par

**Paper:** C. Soss et al., *ScaWL: Scaling k-WL (Weisfeiler–Lehman) Algorithms in Distributed-Memory.* OSTI 2587225.

\par\medskip\textbf{\# Status: ✅ Complete — Score 9/10 (lifted 2026-04-28: added distributed C++/MPI 3-WL run on JLSE chiatta00)}\par

Independent Python/NumPy re-implementation of 2-WL color refinement with
`multiprocessing`-based parallelism; strong-scaling benchmark reproduces the
qualitative shape of the paper's Table 4 on a 10-core workstation.

\par\medskip\textbf{\# Headline numbers (n=200, d=4 random-regular graph, CherryRd)}\par

| p  | time (s) | speedup (ours) | speedup (paper avg) |
|----|---------:|---------------:|--------------------:|
| 1  | 5.82     | 1.00           | 1.00   |
| 2  | 3.08     | 1.89           | 2.38   |
| 4  | 1.68     | 3.46           | 4.26   |
| 8  | 0.94     | 6.17           | 7.64   |
| 16 | 0.79     | 7.41           | 13.20* |

*paper at 16 cores uses 2 physical sockets; we only have 10 physical cores.

\par\medskip\textbf{\# Correctness (all pass)}\par
\par$\bullet$\ Iso pair (random 4-regular, n=30 with relabeling) ✓
\par$\bullet$\ Non-iso 2-regular pair C12 vs 2×C6 (distinguished by 2-WL but not 1-WL) ✓
\par$\bullet$\ Strongly-regular pair (Shrikhande vs 4×4 rook) — match expected (2-WL cannot distinguish) ✓

\par\medskip\textbf{\# Layout}\par
[code block omitted]

\par\medskip\textbf{\# Reproduce}\par
[code block omitted]
\normalsize

The replication was driven by an OpenClaw agent loop (Claude Opus / GPT-5 class models) issuing shell, Python, and (where applicable) GPU jobs. Compute usage and wall-time for individual papers are summarised in the meta-paper's resource table; not separately recorded for this paper unless noted in the Tier-Lift artefact. Sub-agents were spawned for replication-plan generation and (for papers in batches 1--4) for tier-lift extension runs.

\section{What was reproduced}
\begin{itemize}
\item Not recorded.
\end{itemize}

\section{What was missing or incomplete}
\begin{itemize}
\item Not recorded.
\end{itemize}
The dominant blocker categories for this paper, mapped onto the meta-paper's 9-category friction taxonomy (\S4), are listed in \S\ref{sec:friction} below.

\section{Quantitative agreement}
Per-quantity numerical comparisons are not separately tabulated in the structured evaluation record for this paper beyond the agreement rationale below; where the rationale references specific numbers, they are reproduced verbatim. A full numerical comparison table is recorded in the per-replication artefacts (see \S\ref{sec:repro}). For papers below 8/8, the agreement rationale identifies the specific quantities that diverge.

\paragraph{Agreement rationale.} Not recorded.

\section{Score and friction-point tagging}\label{sec:friction}
\begin{itemize}
\item \textbf{Coverage:} 8/10 --- 2-WL re-implementation + correctness oracles + strong-scaling reproduction (single node); plus distributed 3-WL C++/MPI run on chiatta00 with strong-scaling sweep at $n\in\{30,60,80,100\}$, $P\in\{1{,}2{,}4{,}8{,}16{,}32{,}64\}$. Multi-node InfiniBand runs and full UFL benchmark suite still out of scope.
\item \textbf{Agreement:} 9/10 --- Qualitative shape of the strong-scaling curve reproduced for both 2-WL and 3-WL on the available single-node hardware. End-to-end $26.4\times$ speedup at $P{=}8\times\mathrm{OMP}{=}16$ on chiatta00 ($n{=}100$). Paper's headline $2{,}193\times$ multi-node figure not attempted (out of scope on one node). Memory feasibility of 3-WL confirmed: $\approx 8$~MB at $n{=}100$, refuting the prior report's $\sim$100~GB Python-overhead estimate by four orders of magnitude.
\end{itemize}

\noindent\textbf{Friction points encountered (taxonomy tags):}
\begin{itemize}
\item \textbf{F5}: Force-field / hyperparameter drift (unspecified configuration)
\end{itemize}

\section{Follow-on questions}
\begin{itemize}
\item Not recorded.
\end{itemize}

\section{Reproducibility pointers}\label{sec:repro}
\begin{itemize}
\item \textbf{Paper directory:} \texttt{2587225-ScaWL-Scaling-k-WL-Weisfeiler-Lehman-Algorithms-in}
\item \textbf{Replication artefacts:}
\begin{itemize}
\item \texttt{/Users/stevens/Dropbox/REPLICATE-PROJECT/2587225-ScaWL-Scaling-k-WL-Weisfeiler-Lehman-Algorithms-in/replication}
\end{itemize}
\item \textbf{Replication-dir hint from JSONL:} \texttt{/Users/stevens/Dropbox/REPLICATE-PROJECT/2587225-ScaWL-Scaling-k-WL-Weisfeiler-Lehman-Algorithms-in}
\item \textbf{Status:} tier\_lift\_v2\_2026\_04\_27
\end{itemize}

To re-run, change directory into the paper directory above and consult its \texttt{README.md} for the entry-point script. Most replications follow the convention \texttt{replication/run\_*.sh} or \texttt{replication/<subdir>/run.py}.

\end{document}
